\chapter{10}

\section{Tessin (CH)}

«Gio-Gio, hallo-o,» die Worte kamen von weit her, «aufwachen,
Kaffeezeit.»

Giovanni tätschelte Louis’ Schulter.

Louis öffnete langsam seine Augen und blickte zu Giovanni hinüber.

Sogleich realisierte er, dass sie die Grenze bereits passiert hatten.
Giovanni musste die Zollkontrolle rasch hinter sich gebracht haben. Und
er hatte sich nicht einmal ausweisen müssen?

«Danke, Giovanni.» Die Worte fielen aus ihm heraus.

«Wofür denn, junger Mann? Ich freue mich über deine Gesellschaft,» und
Giovannis faltiges Gesicht leuchtete, «ich fahre sonst stundenlang
alleine.»

Giovannis Fahrt würde im Tessin enden. Anschliessend werde er wieder
nach Mailand zurückkehren, erklärte er.

Er schlug Louis vor, bei seinen Trucker-Freunden nachzufragen, wer
direkt nach Zürich fuhr.

Es ging störungsfrei voran. Giovanni schickte über Funk Louis’ Anliegen
an eine Reihe seiner LKW-Kollegen. Innerhalb weniger Minuten meldete
sich eine Männerstimme. Die beiden sprachen auffallend vertraut
miteinander. Der Mann bot seine Unterstützung an. Er klang, als wäre
nichts selbstverständlicher, als Louis abzuholen.

Nachdem sie ihren Kaffee getrunken hatten, tauchte bereits ein riesiger
LKW auf, hellgrün, und hielt in ihrer Nähe. Louis war, als würde er
gleich in eine eingeschworene Gemeinschaft aufgenommen.

Ein glatzköpfiger Mann mit freundlichem Blick kam auf sie zu. Giovanni
begrüsste ihn mit kumpelhafter Herzlichkeit und wiederholte in wenigen
Worten seine Bitte. Dabei zeigte er auf Louis.

Louis war beeindruckt, wie unkompliziert seine Reise weiterging.
Gleichermassen irritierte ihn das Tempo. Bis dahin passte alles. Schlag
auf Schlag.

Es war wichtig, aufmerksam zu bleiben.

Er studierte das Gesicht des Mannes und war beruhigt. Er fand darin
weder offene Fragen noch Zweifel.

Mit «Bin Leonard, hi,» stellte sich der Chauffeur vor. Er klopfte Louis
heftig auf die Schulter. Louis zuckte zusammen. Fing sich wieder und
suchte nach seinem neuen Namen.

«Giovanni,» er bemühte sich zu lächeln, «danke, dass Sie mich
mitnehmen.»

«Kein Thema, für meinen Freund tu ich fast alles.»

Mit diesen Worten drehte sich Leonard zu Giovanni um. Er grinste, beugte
sich zu Giovanni hinunter, hielt ihn über die Schulter und wünschte ihm
eine leichte Fahrt.

Der verschwommene Blick durch die Brille strengte Louis an. Seine Augen
tränten.

Neben ihnen krachten ununterbrochen Personenwagen, LKWs und riesige
Trucks vorbei. Dazwischen zwängte sich grelles Gehupe.

«Steig ein, deine Lady wird sich bestimmt freuen, wenn du auftauchst,»
schrie Leonard Louis zu. Dabei blinzelte er verschwörerisch. Jetzt
musste Louis lachen. Ziemlich herzhaft, seit langem. Was für eine
Inszenierung. Wie grotesk das alles war.

Und wieder dachte er, dass er konzentriert bleiben musste, um seine
Geschichte nicht zu verlieren.

Leonard war Deutschschweizer.

Mit dem Wechsel von Land und Sprache entfernte sich Louis von Mailand
und der damit verbundenen Katastrophe. Für ein paar Momente gelang ihm,
sich in beruhigende Gedanken zu retten. Vielleicht hatte er sich zu sehr
hineingesteigert. Und vielleicht war alles weit weniger dramatisch
ausgegangen?

Leonard war wortkarg. Typ zugeknöpft. Das Radio lief dauernd und viel zu
laut. Zum einen. Zum andern war es gut. Louis stellte sich müde und
unbeteiligt. Liess sich zwischendurch auf kurze Gespräche ein. Über
belanglose Dinge, wie die Hitze dieses Sommers. Oder über seine
Beziehung zu dieser vermeintlichen Freundin. Dann stellte er sich
schlafend. Und Leonard liess ihn in Ruhe.

Bis Louis die Augen öffnete. Die Ortsschilder an den Autobahnausfahrten
zeigten, dass sie gut vorangekommen waren. Zürich war nah. Leonard
schaute kurz zu ihm hinüber.

Im selben Moment raste Adrenalin durch Louis’ Körper. Unerwartet hatte
Leonard ihm die Frage gestellt.

«Wo warst du eigentlich in Italien?»

Louis glaubte in seiner Stimme Unsicherheit wahrzunehmen, oder Skepsis?
Oder beides? Sein Kopf begann zu schmerzen. Doch dann ergaben sich die
Sätze ganz von alleine.

«Ich war in der Toscana bei meinem Onkel, habe ihm in seinem Restaurant
geholfen, so eine Art Auszeit für mich. Ich wollte mir klar werden, was
ich eigentlich will. Meine Freundin spricht von Familie, Kindern, du
weisst schon.»

Leonard wackelte mit dem Kopf hin und her. Rollte die Augen und lachte.

«Ja, ja, die Frauen. Ich versteh dich. Sehr, sehr kompliziert, das kenn
ich. Mein Gott.»

Er winkte ab und zündete sich eine Zigarette an.

Damit war das Gespräch beendet.

Leonard bog auf die nächste Raststätte ein, um zu tanken.

Louis begleitete ihn zur Kasse.

«Was schuld ich dir?»

«Einen Hunderter, hahaha,» Leonard lachte laut und verzog seinen Mund,
«nix natürlich. Ehrensache. Giovanni hat mich darum gebeten.»

Als Leonard an der Kasse zu Zigaretten griff, nahm sie Louis ihm sanft
aus der Hand und bezahlte. Leonard nickte stumm.

Louis setzte sich ein letztes Mal neben ihn in die Führerkabine.

Das Schild «Autobahnausfahrt Zürich Süd» tauchte auf. Es passte. Er bat
Leonard, rauszufahren und ihn abzusetzen.

«Ich lass dich beim Rotlicht raus, schau, dort vorne, ok? Hier kann ich
schlecht anhalten.»

Leonard konzentrierte sich mit Blick in den Rückspiegel auf den Verkehr.
